\section{Определение однолистной функции, конформного отображения.
Свойства конформных
отображений.}\label{ux43eux43fux440ux435ux434ux435ux43bux435ux43dux438ux435-ux43eux434ux43dux43eux43bux438ux441ux442ux43dux43eux439-ux444ux443ux43dux43aux446ux438ux438-ux43aux43eux43dux444ux43eux440ux43cux43dux43eux433ux43e-ux43eux442ux43eux431ux440ux430ux436ux435ux43dux438ux44f.-ux441ux432ux43eux439ux441ux442ux432ux430-ux43aux43eux43dux444ux43eux440ux43cux43dux44bux445-ux43eux442ux43eux431ux440ux430ux436ux435ux43dux438ux439.}

\textbf{Определение:} \(w=f(z)\) \emph{однолистная} на \(G\), если
\(\forall z_{1},z_{2}\in G, \ z_{1}\neq z_{2}\implies f(z_{1})\neq f(z_{2})\)

\textbf{Определение:} Отображение, осуществляемое функцией
\(w=f(z), z\in G\), называется \emph{конформным}, если удовлетворяет
условиям:

\begin{enumerate}
\def\labelenumi{\arabic{enumi}.}
\tightlist
\item
  \(f(z)\) аналитическая в \(G\);
\item
  \(f(z)\) однолистная функция;
\item
  \(f'(z)\neq 0\quad\forall z\in G\)
\end{enumerate}

\textbf{Свойства конформных отображений:}

\begin{enumerate}
\def\labelenumi{\arabic{enumi}.}
\tightlist
\item
  При конформным отображении \(w=f(z)\) образом области \(G\) является
  область расширенной комплексной области плоскости \(w\).
\item
  Отображение, обратное к конформному отображению, также является
  конформным.
\item
  Суперпозиция (последовательное выполнение) конформных отображений
  также является конформным
\item
  При конформном отображении \(w=f(z)\) области \(G\) в каждой конечной
  точке \(z_{0}\in G\) \emph{линейное растяжение} одинаково для всех
  гладких кривых с началом в точке \(z_{0}\) и равно \(|f'(z_{0})|\),
  т.е. \emph{образом окружности} \(|z-z_{0}|=\rho\) с точностью до
  \(o(\rho)\) является \emph{окружность} \(|w-w_{0}|=\rho |f'(z_{0})|\),
  где \(w_{0}=f(z_{0})\)
\item
  При конформном отображении \(w=f(z)\) области \(G\) углы между кривыми
  с началом в точке \(z_{0}\in G, z_{0}\neq \infty\), или с концом в
  точке \(z_{0}=\infty\) равен углу между образами этих кривых в точке
  \(w_{0}=f(z_{0})\); при этом если \(z_{0}\neq \infty\), то угол
  поворота всех кривых с началом в точке \(z_{0}\) равен
  \(\arg f'(z_{0})\).
\end{enumerate}
